\section{Preliminaries}
\label{sec:prelims}

\subsection{Notation}
\label{sec:notation}

\paragraph{Strings.}
Write $\bits$ for the binary alphabet. For an integer $n \geq 0$, $\bits^n$
denotes the set of $n$-bit strings, $\bits^* = \bigcup_{n \geq 0} \bits^n$
the set of all finite-length bit strings, and $\varepsilon$ the empty
string. For $x \in \bits^*$, $|x|$ is the bit length and $x \concat y$ is
the concatenation; the bitwise XOR of two equal-length strings $x, y \in
\bits^n$ is $x \oplus y$. A \emph{byte string} is an element of
$(\bits^8)^*$; its byte length is denoted $\|x\| = |x|/8$. For an
integer $0 \leq n < 2^{64}$, $\langle n \rangle_8 \in \bits^{64}$ is its
eight-byte little-endian encoding. For a finite or infinite bit string
$x$ and $0 \leq \tau \leq |x|$ when $x$ is finite,
$\lfloor x \rfloor_\tau \in \bits^\tau$ denotes the leading $\tau$ bits
of $x$.

\paragraph{Sampling and probability.}
We write $x \sample S$ for sampling $x$ uniformly at random from a finite
set $S$, with implicit coins independent across samples. For an event $E$
in a probability experiment, $\Pr[E]$ denotes its probability.

\paragraph{Sets of maps.}
For a set $S$, $\Perm(S)$ denotes the set of permutations on $S$; for a
positive integer $b$, $\Perm(\bits^b)$ is the set of $b$-bit
permutations.

\paragraph{Games and advantages.}
We write $\mathcal{A}^{O_1, \dots, O_k}$ for a (possibly randomized)
adversary with oracle access to $O_1, \dots, O_k$, and
$\Pr[\mathsf{G}^{\mathcal{A}} \Rightarrow 1]$ for the probability that a
game $\mathsf{G}$---an experiment that runs $\mathcal{A}$ against
specified oracles and returns a Boolean outcome---returns $1$. For a
distinguishing experiment between two games $\mathsf{G}_0, \mathsf{G}_1$,
the advantage of $\mathcal{A}$ against a scheme $\Pi$ in notion
$\mathsf{X}$ is the absolute difference
\[
  \Adv^{\mathsf{X}}_\Pi(\mathcal{A})
  = \bigl| \Pr[\mathsf{G}_0^{\mathcal{A}} \Rightarrow 1]
         - \Pr[\mathsf{G}_1^{\mathcal{A}} \Rightarrow 1] \bigr|.
\]
For a win-event experiment without a hidden challenge bit it is instead
$\Pr[\mathsf{G}^{\mathcal{A}} \text{ wins}]$, the probability the game
returns $1$; a bit-guessing experiment with a hidden challenge bit uses
the rescaled form
$\bigl|2 \Pr[\mathsf{G}^{\mathcal{A}} \text{ wins}] - 1\bigr|$, as in the
multi-user notion of \Cref{sec:mu-ind}.

\paragraph{Concrete security.}
All bounds in this paper are concrete: each advantage is an explicit
function of the resource caps introduced in
\Cref{sec:resources}. We make no use of asymptotic or negligibility
conventions.

\paragraph{Reference summary.}
\Cref{tab:notation} summarizes the resource caps and loss terms used
throughout. The caps are defined formally in \Cref{sec:resources}, the
loss terms in \Cref{sec:loss-terms}.

\begin{table}[ht]
  \centering
  \small
  \begin{tabular}{l l l}
    \toprule
    Symbol & Meaning & Defined \\
    \midrule
    $N$        & Construction interface cost (\Keccakp{} calls)       & \cref{sec:resources} \\
    $\Nprim$   & Direct primitive-interface queries                   & \cref{sec:resources} \\
    $\Ntot$    & $N + \Nprim$ (total query-sequence cost)             & \cref{sec:resources} \\
    $\qv$, $\Qv$ & Valid verification queries / cap                   & \cref{sec:resources} \\
    $\qRO$, $\QRO$ & Direct random oracle queries / cap (RO-model)    & \cref{sec:resources} \\
    $\nleaf$, $\Nleaf$ & Distinct construction-generated final leaf transcripts / cap & \cref{sec:resources} \\
    \midrule
    $\Lift_c(\Ntot)$        & \cite{BDPV08} flat sponge loss   & \cref{sec:loss-terms} \\
    $\KeyGuess_k(Q)$        & Ideal system key guess           & \cref{sec:loss-terms} \\
    $\RootColl_\tau(Q, s)$  & $s$-way root-collision           & \cref{sec:loss-terms} \\
    $\RootPre_\tau(Q)$      & Root-preimage                    & \cref{sec:loss-terms} \\
    $\LeafColl_\tau(Q)$     & Known-key leaf-collision         & \cref{sec:loss-terms} \\
    \bottomrule
  \end{tabular}
  \caption{Resource caps and loss terms.}
  \label{tab:notation}
\end{table}

\subsection{Sponges and Duplexes over \texorpdfstring{\Keccak{}}{Keccak-p}}
\label{sec:sponge}

\paragraph{The \Keccak{} permutation.}
The \Keccak{} family of permutations, defined in
FIPS~202~\cite{FIPS202} as the generalization of the Keccak design's
Keccak-$f$ permutations to a free round count, is parameterized
by a state width $b$ and a round count $n_r$; we write $\Keccak[b, n_r]
\in \Perm(\bits^b)$ for the corresponding permutation. \TW{} uses
\Keccakp{}, the same primitive as
KangarooTwelve~\cite{BDPVVV16}, fixed throughout. In the
public permutation security model of \Cref{sec:games-pp} and the
flat sponge lift of \Cref{app:flat-sponge-lift}, this permutation is
modeled as a uniformly random element of $\Perm(\bits^{1600})$.

\paragraph{The sponge construction.}
Following \cite{BDPV08}, fix a permutation $\pi \in \Perm(\bits^b)$, a
\emph{sponge-compliant} padding rule $\mathsf{pad}$, and a rate $1 \leq
r < b$ with capacity $c = b - r$. A padding rule appends
$\mathsf{pad}[r](|M|)$, a bit string determined by the input length and
the rate, to an input $M$, extending it to a sequence of $r$-bit
blocks; it is sponge-compliant
\cite[Definition~1]{BDPV11} if $M \concat \mathsf{pad}[r](|M|)$ is never
empty and, for all $M \neq M'$ and all $n \geq 0$, $M \concat
\mathsf{pad}[r](|M|) \neq M' \concat \mathsf{pad}[r](|M'|) \concat
0^{n r}$. The sponge function $\Sponge[\pi, \mathsf{pad}, r]
: \bits^* \times \mathbb{N} \to \bits^*$ maps an input $M$ and an output
length $\ell$ to an $\ell$-bit string by absorbing $M \concat
\mathsf{pad}[r](|M|)$ into the all-zero $b$-bit state in $r$-bit blocks
(each separated by an application of $\pi$) and squeezing $\ell$ bits in
$r$-bit blocks from the resulting state. \cite[Theorem~2]{BDPV08} shows that
$\Sponge[\pi, \padten, r]$, where the simple rule $\padten$ appends a
single $1$ bit followed by the minimum number of $0$ bits reaching a
multiple of $r$, is indifferentiable from a random oracle when
$\pi$ is uniformly random. We bound its distinguishing advantage by
$\Lift_c(\Ntot) = \Ntot(\Ntot+1) / 2^{c+1}$, derived in
\Cref{lem:lift-bound} from the product upper bounds of
\cite[Lemmas~4 and~5]{BDPV08}.

\paragraph{The duplex construction.}
The \emph{duplex} construction of \cite{BDPV11} extends the sponge to
a stateful interface. A duplex object $\Duplex[\pi, \mathsf{pad}, r]$
is initialized to the all-zero state. Each call $\duplexing(\sigma,
\ell)$, with $\ell \leq r$, absorbs an input string $\sigma$ of at most
$\rho_{\max}(\mathsf{pad}, r)$ bits via $\pi$ and returns the first
$\ell$ bits of the resulting state. Here $\rho_{\max}(\mathsf{pad}, r)
= \min\{x : x + |\mathsf{pad}[r](x)| > r\} - 1$ is the largest input
length for which the padded input fits in a single $r$-bit block; for
the $\padtenone$ rule, $\rho_{\max}(\padtenone, r) = r - 2$
\cite{BDPV11}. \cite[Lemma~3]{BDPV11}
establishes the duplex-to-sponge equality: for any call
sequence $((\sigma_0, \ell_0), \dots, (\sigma_i, \ell_i))$, the $i$-th
output equals
\[
  \Sponge[\pi, \mathsf{pad}, r]\bigl(
    \sigma_0 \concat \mathsf{pad}_0 \concat \cdots \concat \sigma_i,\,
    \ell_i\bigr),
\]
where $\mathsf{pad}_j = \mathsf{pad}[r](|\sigma_j|)$ is the padding
the duplex appends in its $j$-th call; the flattened input carries
these interior paddings explicitly, and the sponge appends only the
final $\mathsf{pad}_i$. The flattening map is
injective by \cite[Lemma~4]{BDPV11}: the duplex input-string sequence is
recoverable from the flattened sponge input.

\paragraph{\TW{} instantiation.}
\TW{} fixes the padding rule $\padtenone{}$ and the rate $r = 1344$, so
the capacity is $c = 256$ bits. The largest input length per duplex call
is therefore $\rho_{\max}(\padtenone, 1344) = 1342$ bits.

\paragraph{Bridge to the \texorpdfstring{$H$}{H}-input domain.}
The indifferentiability theorem of \cite{BDPV08} is stated for the
simpler $\padten{}$ padding over its unpadded $H$-input domain, while
\TW{} uses $\padtenone{}$ internally. \cite[Theorem~3]{BDPV11} supplies
an injective preprocessing bridge $\bridge[r, r_{\max}]$ between the two
domains. For \TW{}, $r = r_{\max} = 1344$, so the bridge reduces to
appending $1 \concat 0^q$ with $q$ determined by the input length modulo
$r_{\max}$; on the bridge image, $q$ is recovered as the trailing-zero
count. \Cref{sec:ro-model} instantiates this bridge on typed \TW{} duplex
transcript prefixes as the map $\flatmap(\cdot)$ used by the
random oracle proofs throughout the paper.

\subsection{Authenticated Encryption with Associated Data}
\label{sec:ae-defs}

\paragraph{Syntax.}
A \emph{nonce-based authenticated encryption scheme with associated
data} (AEAD scheme) is a pair $\SE = (\Enc, \Dec)$ of deterministic
algorithms with associated key space $\mathcal{K} \subseteq (\bits^8)^*$,
nonce space $\mathcal{U} \subseteq (\bits^8)^*$, associated data space
$\mathcal{A} \subseteq (\bits^8)^*$, message space $\mathcal{M} \subseteq
(\bits^8)^*$, and ciphertext space $\mathcal{C} \subseteq (\bits^8)^*$. The
encryption algorithm $\Enc : \mathcal{K} \times \mathcal{U} \times
\mathcal{A} \times \mathcal{M} \to \mathcal{C}$ produces a ciphertext
$C = \Enc_K(U, A, M)$. The decryption algorithm $\Dec : \mathcal{K}
\times \mathcal{U} \times \mathcal{A} \times \mathcal{C} \to
\mathcal{M} \cup \{\bot\}$ returns either a plaintext or the
rejection symbol $\bot$. The \emph{tag length in bytes} is the constant
$\|\Enc_K(U, A, M)\| - \|M\|$, fixed by the scheme.
We write nonce values as $U$ to keep $N$ available for the construction-cost
resource variable used in the concrete bounds.
In this generic syntax, $C$ denotes the full AEAD ciphertext, including
any tag material. The \TW{} algorithm section (\Cref{sec:tw128}) writes this same value as
$C \concat T$, with $C$ the ciphertext body and $T$ the root tag.

\paragraph{Correctness.}
For every $(K, U, A, M) \in \mathcal{K} \times \mathcal{U} \times
\mathcal{A} \times \mathcal{M}$, $\Dec_K(U, A, \Enc_K(U, A, M)) = M$.
An AEAD scheme is \emph{tidy} if every accepting decryption is consistent
with encryption: whenever $\Dec_K(U, A, C) = M \neq \bot$, one has
$\Enc_K(U, A, M) = C$.

\paragraph{All-in-one AE.}
We use the all-in-one authenticated encryption notion of \cite{RS06},
instantiated for nonce-based AE following \cite{BT16}. The experiment
grants either the real pair $(\Enc_K,\Ver_K)$ or the ideal pair
$(\Rand,\Replay)$, where $\Rand$ returns a fresh uniform byte string of
the ciphertext length on each encryption query and $\Replay$ accepts
exactly the tuples $\Rand$ produced. The advantage
$\Advae_\SE(\mathcal{A})$ is the distinguishing advantage between these
two worlds over nonce-respecting adversaries. This all-in-one notion is
instantiated for \TW{} in \Cref{sec:games-pp}.

\paragraph{Integrity under release of unverified plaintext.}
For the RUP setting of \cite{ABLMMY14}, a conventional decryption algorithm is
separated into a plaintext-computing algorithm and a verification algorithm.
The integrity notion \textsf{INT-RUP} gives the adversary access to encryption,
plaintext computation, and verification, and asks whether verification accepts a
fresh ciphertext and tag not returned by encryption. As in the nonce-IV setting
of \cite{ABLMMY14}, encryption queries are nonce-respecting, while plaintext
computation and verification may be queried on arbitrary nonces. This notion
protects only ciphertext integrity; privacy under release of unverified
plaintext requires plaintext awareness in addition to ordinary encryption
privacy.

\subsection{Multi-User Indistinguishability}
\label{sec:mu-ind}

The multi-user indistinguishability notion of \cite[Fig.~3]{BT16}
extends the same real-vs-ideal oracle pair to an adversary-selected set
of users. The number of users $u$ is adversary-determined: it is the
number of $\mathsf{New}$ queries the adversary makes, i.e., the final
value of the user counter $v$ in the game below. A challenge bit $b
\sample \bits$ selects between the
real construction interface (real ciphertexts, real verification) and
an ideal interface (uniform strings of matching length, verification
that accepts only previously-issued encryption tuples). The adversary
also accesses a public primitive interface: the real permutation
$(\pi,\pi^{-1})$ in the real world and the simulator
$(\Sim^{\ROflat},\Sim^{-1,\ROflat})$ in the ideal world; this
world-dependent primitive interface is our modification---\cite{BT16}
work in the ideal-cipher model and give both worlds the same
ideal-cipher oracles---using the same ideal-side simulator convention
as the single-user AE-style games. The adversary
outputs a guess $b'$ and wins if $b' = b$.

\paragraph{Game and advantage.}
The game $\mathsf{G}^{\mathsf{mu\text{-}ind}}_\SE(\mathcal{A})$ of
\cite[Fig.~3]{BT16} samples $b \sample \bits$ and maintains a user
counter $v$, a set $V$ of $(d, U)$ pairs already used by encryption,
and a set $W$ of recorded encryption tuples; its $\mathsf{New}$,
$\mathsf{Enc}$, and $\mathsf{Vf}$ oracles are instantiated for \TW{} in
\Cref{sec:mu-setting}, together with the added public
primitive interface---the real permutation when $b = 1$ and the
simulator when $b = 0$. We write the user index as $d$ and the
associated data input as $A$, where \cite{BT16} writes $i$ and $H$, and
take the absolute value of their signed advantage so that
\[
  \begin{aligned}
  \Advmuind_\SE(\mathcal{A})
  &= \left|2 \Pr[\mathsf{G}^{\mathsf{mu\text{-}ind}}_\SE(\mathcal{A})
          \text{ returns } 1] - 1\right| \\
  &= \bigl|
      \Pr[\mathcal{A}^{\mathsf{Enc}, \mathsf{Vf}, b=1} \Rightarrow 1]
      - \Pr[\mathcal{A}^{\mathsf{Enc}, \mathsf{Vf}, b=0} \Rightarrow 1]
    \bigr|
  \end{aligned}
\]
the second equality holding because $b$ is uniform. Thus
$\Advmuind_\SE$ is the two-game distinguishing advantage between the
real $(\Enc, \Ver)$ pair and the ideal $(\Rand, \Replay)$ pair
instantiated per user. The single-user case $u = 1$ recovers the
all-in-one AE notion of \Cref{sec:ae-defs}, augmented with the public
primitive interface (the real permutation in the real world and the
simulator in the ideal world).

\paragraph{Per-user nonce uniqueness.}
Nonce uniqueness is enforced per user via the set $V$: the same nonce
may appear under distinct users but not twice under the same user.
This is strictly weaker than the global nonce-uniqueness requirement
that would force each nonce value $U$ to appear under at most one user.

\subsection{Committing Notions and Root Tag Hash Security}
\label{sec:commit-defs}

\cite{BH22} formalize a family of committing security notions for
nonce-based AEAD schemes, parameterized by how much of the
encryption input is committed and by whether commitment is
established through decryption (the D-notions) or through encryption
agreement (the E-notions). We recall their definitions and the
$s$-way generalization, and then frame the root tag itself as a hash function
whose \cite{RS04} security---collision, second-preimage, and preimage
resistance---captures the binding and preimage properties expected of a compact
commitment. The AEAD committing notions are then proved from the same tag
analysis in \Cref{sec:root-tag-hash-security}.

\paragraph{What-is-committed functions.}
For $\ell \in \{1, 3, 4\}$, the \emph{what-is-committed} function
$\WiC_\ell$ projects an encryption input $(K, U, A, M) \in \mathcal{K}
\times \mathcal{U} \times \mathcal{A} \times \mathcal{M}$ to the
portion it should commit:
\[
\begin{aligned}
  \WiC_1(K, U, A, M) &= K, \\
  \WiC_3(K, U, A, M) &= (K, U, A), \\
  \WiC_4(K, U, A, M) &= (K, U, A, M).
\end{aligned}
\]

\paragraph{CMTDs-$\ell$ and CMTs-$\ell$ games.}
For $s \geq 2$, \cite[Fig.~5]{BH22} defines the $s$-way committing
games on tuples $(K_i, U_i, A_i, M_i)_{i=1}^s$ with pairwise distinct
$\WiC_\ell$-images, all entries in the scheme's syntactic domain and
none equal to $\bot$. The $s$-way CMTDs-$\ell$ (decryption) game has
the adversary output a ciphertext $C$ alongside the tuples and returns
$1$ if $\Dec_{K_i}(U_i, A_i, C) = M_i$ for every $i$; the printed
return line of \cite[Fig.~5]{BH22} writes $C_i$, but the game outputs a
single ciphertext $C$. The $s$-way CMTs-$\ell$ (encryption) game has the
adversary output the tuples alone and returns $1$ if all encryptions
$\Enc_{K_i}(U_i, A_i, M_i)$ are equal. Advantages $\Advcmtds{\ell}_\SE(\mathcal{A})$ and
$\Advcmts{\ell}_\SE(\mathcal{A})$ are the probabilities that the
adversary wins the respective game; the parameter $s$ is suppressed
from the notation when clear from context. The two-tuple case $s = 2$
is the original CMTD-$\ell$/CMT-$\ell$ pair of \cite[Fig.~4]{BH22}.
These definitions are instantiated for \TW{} in \Cref{sec:games-pp},
with the public permutation lift deferred to that section.

\begin{definition}[Root tag wrapper]
\label{def:htw}
For every \TW{} encryption input $X = (K,U,A,M)$ in the scheme domain,
define the root tag wrapper by
\[
  \HTW(K,U,A,M) = T
  \quad\text{where}\quad
  \Enc_K(U,A,M) = C \concat T.
\]
We also write $\HTW(X)$ for $\HTW(K,U,A,M)$. Equivalently, \HTW{} runs
\TW{} encryption and discards the ciphertext body.
\end{definition}

\paragraph{The root tag as a hash.}
Many uses of commitment store, transmit, or compare a short tag rather
than the full ciphertext, and ask that the short string alone be
binding---the compactly committing AEAD goal introduced by \cite{GLR17}
for message franking. \TW{} ciphertexts are written $C \concat T$ with
$T$ the $\tauroot$-bit root tag, and we capture this goal directly by
treating \Cref{def:htw} as a hash function. Keyed by the permutation,
this is a hash family, and we ask that it satisfy the hash-function
security notions of \cite{RS04}, over the \TW{} input domain, together
with the local $s$-way multicollision extension used for the committing
theorems. For fixed byte lengths $a$ and $\mu$, write
\[
  \mathcal{X}_{a,\mu}
  =
  \bits^k \times \bits^\nu \times (\bits^8)^a \times (\bits^8)^\mu
\]
for the fixed-length slice of valid inputs, with bit length
$m_{a,\mu} = k + \nu + 8a + 8\mu$.
\begin{description}
  \item[Collision resistance (Coll).] No adversary finds distinct inputs
    $X_1, \dots, X_s$ (any $s \geq 2$) with
    $\HTW(X_1) = \cdots = \HTW(X_s)$. The case $s=2$ is the RS04 Coll
    notion; larger $s$ is the corresponding multicollision extension.
  \item[Preimage resistance (Pre).] For fixed $a,\mu$, given
    $\HTW(X^*)$ for a uniform $X^* \sample \mathcal{X}_{a,\mu}$, no
    adversary finds any valid input $X$ with $\HTW(X) = \HTW(X^*)$.
  \item[Second-preimage resistance (Sec/eSec).] The standard and
    everywhere forms of second-preimage resistance of \cite{RS04}, on
    each fixed input slice $\mathcal{X}_{a,\mu}$, follow from
    collision resistance by \cite[Proposition~6]{RS04}.
  \item[Everywhere preimage resistance (ePre).] For a target tag $T^*$ fixed in
    advance, no adversary finds $X$ with $\HTW(X) = T^*$.
\end{description}
The advantages $\AdvColl_\TW$, $\AdvPreSlice{a}{\mu}_\TW$, and
$\AdvePre_\TW$ are the respective success probabilities, instantiated for
\TW{} in \Cref{sec:games-pp}. The Pre guarantee is proved directly in
\Cref{thm:pre}, while the second-preimage guarantees are inherited from Coll
in \Cref{cor:second-preimage-rs04}. Coll is the binding property of the tag
taken as a compact commitment to $(K, U, A, M)$, Pre is the random-opening
preimage property on fixed input slices, and ePre is the fixed-target
preimage property for query-independent tags.
The \cite{BH22} committing notions are the collision-facing AEAD
projection of the same tag analysis: a winning opening gives distinct
inputs sharing a tag, while the dedicated CMTDs-4 proof uses the
all-bodies-equal game structure to avoid the leaf-collision slack of the
general hash collision theorem (\Cref{sec:root-tag-hash-security}).

\paragraph{Relations.}
\cite[Figs.~4 and~5]{BH22} record the implications among the notions.
For the D-notions we use:
\[
  \mathsf{CMTDs}\text{-}4 \;\Rightarrow\;
  \mathsf{CMTDs}\text{-}\ell \text{ for } \ell \in \{1, 3\}
\]
(both immediate: pairwise distinct keys ($\ell = 1$) or pairwise
distinct $(K, U, A)$ triples ($\ell = 3$) are \emph{a fortiori} pairwise
distinct full tuples, so every CMTDs-$\ell$ winner is a CMTDs-4 winner).
The E-notions are bounded directly in \Cref{cor:committing-combined} by
the same final-root candidate-collision argument, rather than through an
encryption reduction.
